% Pandoc-LaTeX-Vorlage für CV (kompilieren mit xelatex)
\documentclass[10pt,letterpaper]{article}

% Language and font encoding
\usepackage[ngerman]{babel}
\usepackage{fontspec}
\defaultfontfeatures{Ligatures=TeX}
\setmainfont{Inter}[
  Path = assets/fonts/,
  UprightFont = Inter-Regular.otf,
  BoldFont = Inter-Bold.otf,
  ItalicFont = Inter-Italic.otf,
  BoldItalicFont = Inter-BoldItalic.otf
]

% Silbentrennungsregeln
\hyphenation{
  Linux
  Unix
  Engineering
  Per-for-mance-ana-ly-se
  Be-rech-ti-gungs-mo-del-le
  Be-rech-ti-gungs-kon-zep-te
  Lö-sungs-spe-zi-fi-ka-ti-o-nen
  Un-ter-neh-mens-kun-den-ge-schäft
  Au-to-ma-ti-sie-rung
  Da-ten-mo-del-lie-rung
  Ak-zep-tanz-kri-te-ri-en
  Zu-griffs-kon-trol-le
  Kom-pe-tenz-ori-en-tier-tes
  Un-ter-neh-mens-in-for-ma-tik
  Pro-zess-ma-nage-ment
  Be-rufs-ma-tu-ri-tät
}

% Page Geometry (Zero left margin for the sidebar)
\usepackage[left=0cm, right=1.2cm, top=1.2cm, bottom=1.2cm, noheadfoot]{geometry}

% Core Packages
\usepackage[dvipsnames]{xcolor}
\usepackage{tikz}
\usepackage{enumitem}
\usepackage{fontawesome5}
\usepackage[hidelinks]{hyperref}
\usepackage{microtype}
\usepackage{graphicx}
\usepackage{titlesec}
\usepackage{eso-pic}      % Robuster Weg, um auf JEDER Seite denselben Hintergrund zu zeichnen
\usepackage{paracol}      % Zwei Spalten, die ueber mehrere Seiten fliessen koennen
\usepackage{changepage}   % adjustwidth: Innenabstand der Sidebar, seitenumbruchfest
\usepackage{ragged2e}     % RaggedRight mit Silbentrennung gegen zu grosse Wortabstaende in der schmalen Sidebar

% Color Definitions (Minimalist Black with Deep Cyan-Blue Accent)
\definecolor{accent}{HTML}{000000}
\definecolor{brand}{HTML}{0369A1}
\definecolor{sidebarbg}{HTML}{F8FAFC}
\definecolor{textdark}{HTML}{000000}
\definecolor{textmuted}{HTML}{334155}

\setlength{\parindent}{0pt}
\setlength{\parskip}{0pt}
\providecommand{\tightlist}{%
  \setlength{\itemsep}{0pt}\setlength{\parskip}{0pt}}

% Section styling (Pandoc erzeugt \section aus ## im Markdown)
\titleformat{\section}{\large\bfseries\color{accent}}{}{0pt}{}
\titlespacing*{\section}{0pt}{0.4cm}{0.15cm}
\newcommand{\sectionrule}{\vspace{0.15cm}\noindent{\color{accent}\rule{\linewidth}{0.5pt}}\par\vspace{0.25cm}}

\titleformat{\subsection}{\large\bfseries\color{accent}}{}{0pt}{}
\titlespacing*{\subsection}{0pt}{0.25cm}{0.1cm}

% Experience Item Command
\newcommand{\experienceitem}[4]{%
  \vspace{0.25cm}%
  \noindent{\large\bfseries\color{accent} #1} \hfill {\small\color{brand}\bfseries #3} \\
  \noindent{\textbf{\color{textdark}#2}} \hfill {\small\color{textmuted}\textit{#4}}%
  \vspace{0.15cm}%
}

% Education Item Command
\newcommand{\educationitem}[4]{%
  \vspace{0.2cm}%
  \noindent{\bfseries\color{accent} #1} \hfill {\small\color{brand}\bfseries #3} \\
  \noindent{\textbf{\color{textdark}#2}} \hfill {\small\color{textmuted}\textit{#4}}\par
  \vspace{0.15cm}%
}

% Project Item Command
\newcommand{\projectitem}[3]{%
  \vspace{0.2cm}%
  \noindent{\large\bfseries\color{accent} #1} \hfill {\small\color{brand}\bfseries #2}\par
  \vspace{0.05cm}%
  \noindent{\textbf{\color{textdark}#3}}\par
  \vspace{0.08cm}%
}

% ATS-freundliche Icon-Beschriftung
\newcommand{\atsicon}[2]{#2}

% Contact Item Command (Sidebar)
\newcommand{\contactitem}[2]{%
  {\noindent\hangindent=0.5cm\hangafter=1
  \makebox[0.5cm][l]{\small\color{brand}\raisebox{-0.08em}{#1}}%
  {\small\color{textdark}#2}\par}%
  \vspace{0.15cm}
}

% Sidebar Section Header
\newcommand{\sidebarsection}[1]{%
  \vspace{0.4cm}%
  {\large\bfseries\color{accent} #1}\par%
  \vspace{0.12cm}%
  \noindent\textcolor{accent}{\rule{\linewidth}{0.5pt}}\par%
  \vspace{0.22cm}%
}

% Section rule nach jedem \section und \subsection
\let\oldsection\section
\renewcommand{\section}[1]{\oldsection{#1}\sectionrule}
\let\oldsubsection\subsection
\renewcommand{\subsection}[1]{{\color{accent}\oldsubsection{#1}}\sectionrule}

% Sidebar-Hintergrund: wird bei JEDER Seite neu gezeichnet (eso-pic-Hook mit
% einfachem \rule statt TikZ-Overlay -- das ist zuverlaessiger in Kombination
% mit paracol), erscheint also auf jeder Seite -- nicht nur auf der ersten
% (das war der Kernfehler: im Original stand dieser Code einmalig im Fliesstext).
\AddToShipoutPictureBG{%
  \AtPageLowerLeft{%
    \color{sidebarbg}%
    \rule{0.33\paperwidth}{\paperheight}%
  }%
}

% Spaltenaufteilung: 33% Sidebar, 67% Hauptspalte.
\columnratio{0.33}
\setlength{\columnsep}{0pt}

\begin{document}
\pagestyle{empty}

\begin{paracol}{2}

\begin{adjustwidth}{0.8cm}{0.6cm}
    \RaggedRight
    \vspace{0.4cm}

    \begin{center}
$if(photo)$
      \begin{tikzpicture}
        \clip (0,0) circle (1.85cm);
        \node[anchor=center,inner sep=0pt] at (0,0) {\includegraphics[width=3.7cm]{$photo$}};
      \end{tikzpicture}
$else$
      \begin{tikzpicture}[scale=0.8]
        \draw[draw=brand, fill=white, line width=1.5pt] (0,0) circle (1.2cm);
        \clip (0,0) circle (1.2cm);
        \fill[brand!10] (-1.2,-1.2) rectangle (1.2,1.2);
        \fill[brand!40] (0,0.15) circle (0.38cm);
        \fill[brand!40] (0,-0.9) circle (0.85cm);
      \end{tikzpicture}
$endif$
    \end{center}
    \vspace{0.25cm}

    \begin{center}
    {\LARGE\bfseries\color{accent} $firstname$\\[0.12cm]$lastname$} \\[0.2cm]
    {\small\bfseries\color{brand} $title$}
    \end{center}

    \vspace{0.35cm}
    \noindent{\color{accent}\rule{\linewidth}{0.8pt}}\par
    \vspace{0.35cm}

$if(email)$
    \contactitem{\atsicon{E-Mail}{\faEnvelope}}{\href{mailto:$email$}{$email$}}
$endif$
$if(phone)$
    \contactitem{\atsicon{Telefon}{\faPhone}}{$phone$}
$endif$
$if(address)$
    \contactitem{\atsicon{Standort}{\faMapMarker*}}{$address$}
$endif$
$if(linkedin)$
    \contactitem{\atsicon{LinkedIn}{\faLinkedin}}{\href{https://linkedin.com/in/$linkedin$}{linkedin.com/in/$linkedin$}}
$endif$
$if(github)$
    \contactitem{\atsicon{GitHub}{\faGithub}}{\href{https://github.com/$github$}{github.com/$github$}}
$endif$

$if(skills)$
    \sidebarsection{Kompetenzen}
$for(skills)$
    \textbf{$skills.category$} \\
    {\small $skills.items$}\par
    \vspace{0.25cm}
$endfor$
$if(pagebreak_sidebar)$
\pagebreak
$endif$
$endif$

$if(certificates)$
\nopagebreak
    \sidebarsection{Zertifikate}
$for(certificates)$
    $certificates$\par
$endfor$
$endif$

$if(education)$
\nopagebreak
    \sidebarsection{Aus- und Weiterbildung}
$for(education)$
    \textbf{$education.degree$}\par
    {\small\color{textmuted}$if(education.year)${\color{brand}\bfseries $education.year$} · $endif$$education.institution$}\par
    \vspace{0.18cm}
$endfor$
$endif$

$if(languages)$
\nopagebreak
    \sidebarsection{Sprachen}
$for(languages)$
    \textbf{$languages.name$} ($languages.level$)\par
$endfor$
$endif$

\end{adjustwidth}

\switchcolumn

\begin{adjustwidth}{0.7cm}{0cm}
  \vspace{0.4cm}
  \color{textdark}
  % Flattersatz ohne Silbentrennung: keine gedehnten Wortabstaende
  \hyphenpenalty=10000
  \raggedright

$body$

\end{adjustwidth}

\end{paracol}

\end{document}
